\section{The Resurrection and the Whole of History}

Christ will return in glory, the dead will rise, and all nations will stand before him. The last judgment manifests the truth of each life in the completed history of the world (Mt 25:31--46; CCC 1038--1040). The particular judgment has already fixed the soul's destiny. At the last judgment that destiny reaches its bodily and public consummation. The one who loved or sinned through a body will stand before the Lord as a risen human being. The history in which the person acted will have reached its end.

The body shares in what a person does and suffers. Feeding someone requires food actually given; protecting someone can require bodily effort and danger. The person who died awaiting Christ has performed such acts through the body which will rise. Bodily resurrection therefore completes the judgment of the whole human being. The soul's destiny, already fixed at death, will be fulfilled in the risen person whose earthly acts Christ judges.

Cyril presses this hope upon hearers whose social position might tempt them to imagine themselves overlooked. The poor laborer in the field, the person grinding at the mill, the servant, and the sick person in bed will all be found. Christ does not prefer the landowner to the laborer or the learned to the simple (XV.23). Cyril then names ordinary fidelities: prayer, almsgiving, fasting, chastity, and a marriage faithfully kept. These things are known to the Lord. The Church's expectation of judgment gives eternal seriousness to lives which scarcely enter the public record.

Thomas explains the need for a final judgment by considering all that remains unfinished when an individual dies. A reputation may remain false. Children continue to live. Teaching continues to influence those who receive it. The body can be honored or dishonored. Projects to which a person devoted himself persist or collapse. These continuing relations belong to a life even after its earthly choosing has ended (III, q.~59, a.~5). God will judge them in the completed whole. The final disclosure reveals a truth no observer situated inside history could fully gather.

Thomas is equally explicit that the dead acquire no new merit or demerit through these later effects. Their state as regards the soul is already fixed. The continued history accompanies the reward or punishment; it does not begin another period of personal choice (a.~5, reply~1). This qualification protects the meaning of responsibility. A parent's child remains a free person. A later reader may twist an author's meaning. A good undertaking may be corrupted by successors. God knows precisely what belonged to the deceased's intention, action, negligence, and foreseeable influence, and what belonged to the new choices of other people.

These limits matter when we consider our influence on others. A parent must teach faithfully and repent of harmful example; a child's holiness still requires the child's free response. A teacher must refuse to spread what he knows is false, even though later readers can misuse a truthful explanation. Someone who has injured a reputation must correct the lie where he can, though he may be unable to follow every retelling. Responsibility includes prudent attention to consequences within our power. God knows where that power ended and another person's choice began.

St.~Augustine opens Book~XX of \emph{The City of God} by confessing that Christ will judge the living and the dead. He then dwells on the bewildering contrasts of present life: the innocent condemned, the guilty released, the useful person dying early, the wrongdoer prospering, the just person enduring illness (XX.1--2). These contrasts are familiar to anyone who has watched suffering closely. Augustine does not tell us that prosperity measures virtue or that affliction proves guilt. At the last judgment, the justice of God's judgments will be manifest, including those whose meaning is now hidden from us.

This promise enters the Christian struggle to trust providence. We sometimes demand an explanation before we will consent to fidelity. We want to know why a prayer was not answered as expected, why a cruel person retained influence, or why a good work seemed to fail. Many such questions will remain with us until death. Faith can entrust them to the Lord without calling evil good or ceasing to resist it. His final judgment assures us that the truth is not being lost because we cannot yet see how it will be vindicated.

It also teaches us to distrust the triumph of an apparent success. A reputation maintained by fear may survive for decades. A household may preserve its respectable appearance by making one member bear everyone's blame. A public achievement may depend upon dishonesty hidden from its beneficiaries. Duration and admiration cannot turn these things into justice. Conversely, work which appeared fruitless may have preserved another person's courage, a concealed prayer may have sustained someone in temptation, and an apparently insignificant kindness may have borne abundant fruit. Christ judges the truth of these acts, including what remained hidden from those who performed them.

The last judgment will reveal God's justice and the wisdom of his providence. We can therefore pursue justice while accepting that some wrongs may remain unresolved when we die. A truthful account, protection for someone in danger, and repair of an injury are still worth seeking when we cannot ensure their completion. Where our work ends, Christ's knowledge of those concerned remains complete. We can entrust the unfinished task to him without withdrawing the care we are able to give now.

